\RequirePackage[l2tabu, orthodox]{nag}
\documentclass[11pt, a4paper]{article}

\def\papertitle{Testing Generation of Chess Games Using Transformer Models}
\def\paperauthor{Hoàng Minh Đức}
\def\paperauthoremail{my email}

\title{\papertitle}
\author{\paperauthor \thanks{\paperauthoremail}}
\date{\today}

\usepackage{fontspec}

\setmonofont{Inconsolata}

\usepackage{graphicx}

\usepackage{xurl}
\usepackage{hyperref}
\hypersetup{
    colorlinks=true,
    linkcolor=blue,
    pdftitle={\papertitle},
    breaklinks=true,
    urlcolor=blue,
    citecolor=blue,
}

\usepackage[backend=biber, natbib=true, style=alphabetic, sorting=none]{biblatex}
\addbibresource{main.bib}

\usepackage{setspace}
\onehalfspacing

\usepackage{parskip}

% tabular nonsense
\usepackage{booktabs}
\usepackage{array}

\usepackage{amsmath}

% in case I need to draw stuff
% \usepackage{xskak}
% \usepackage{tikz}

\usepackage[capitalize, noabbrev]{cleveref}
% \usepackage{amssymb}

\begin{document}

\maketitle

\begin{abstract}
    % TODO: add more to this abstract.
    In this paper, we analyze whether decoder-only Transformer models can generate complete (terminated) and valid (no illegal moves or positions) chess games by watching movement of pieces.
\end{abstract}

\section{Introduction}

Computer chess has a very long history, dating back to the earliest days of computer science \citep{silver-2017}. Programs, called \textit{chess engines}, such as Stockfish \citep{thestockfishdevelopers-2026} have achieved superhuman performance, utilizing techniques such as

The Transformer architecture \citep{vaswani-2017} is a very influential architecture, used for the majority of language models. It has seen use in natural language processing, vision, and even analysis of chess itself \citep{ruoss-2024,monroe-2024}. In this paper, we build and analyze Transformer-based generative (decoder-only) models at the ability to synthetically generate chess games.

The

\section{Related Works}

\section{Methodology}

\subsection{Definitions}

A \textit{chess game} is a finite sequence of moves from the standard initial position of chess, in which every move is legal given the current board state. We distinguish two categories:

\begin{itemize}
    \item A \textit{complete chess game} is a chess game that terminates in a decisive condition under the FIDE Laws of Chess \citep{fide-2023}: checkmate, stalemate, insufficient material, the fifty-move rule, or threefold repetition.
    \item An \textit{incomplete chess game} is a chess game that does not meet the above conditions. This includes games ended by resignation or mutual agreement, whose terminal conditions cannot be inferred from the move sequence alone.
\end{itemize}

Moves are represented using Universal Chess Interface (\textit{UCI}) notation, in which each move is defined by picking up a piece in a \textit{starting square}, and placing it in the \textit{destination square}. Even though Standard Algebraic Notation (\textit{SAN}) is more frequently used, UCI notation has a much lower overall token vocabulary, since each token does not need to carry extra information that SAN explicitly represents, e.g. the piece being moved.

\subsection{Data}

For this particular task, we use the Lichess open databases, which consists entirely of games by human and engine players of various ELO levels.

\subsection{Model}

\subsubsection{Model Architecture}

The model is based on the Transformer architecture, itself based on the attention mechanism \citep{vaswani-2017}. In particular, we will use OpenAI's GPT-2 model \citep{radford-2019}, and it's implementation in HuggingFace's Transformer Library.

\section{Results}

\section{Conclusion}

\section{Acknowledgments}

I would like to acknowledge Austin Davis for preprocessing and creating the UCI-notation version of the Lichess monthly databases \citep{davis-2024}.

\appendix

% \pagebreak
\printbibliography
\end{document}
